\documentclass[10pt]{article}
\usepackage[T1]{fontenc}
\usepackage[utf8]{inputenc}
\usepackage[french]{babel}
\usepackage{amsmath,amssymb,mathrsfs,mathpazo}
\usepackage{geometry}
\geometry{a4paper,margin=2cm}
\usepackage{graphicx,array,tabularx,tabularray,colortbl,multirow,multicol,booktabs,enumitem,pifont,ulem,fancybox,diagbox,cancel}
\usepackage{pstricks,pstricks-add,pst-plot,pst-tree,pst-node,pst-text,pst-func,pst-math,pst-bspline,pst-3dplot,pst-eucl,pst-solides3d}
\let\pscalculate\undefined
\usepackage{pst-bezier}
\usepackage{tikz,pgfplots,tkz-tab}
\usetikzlibrary{arrows,arrows.meta,patterns,calc,positioning,trees,decorations.markings,shapes.geometric,shapes.misc,shapes.arrows,intersections,decorations.pathmorphing,decorations.pathreplacing,backgrounds,shadows,babel}
\pgfplotsset{compat=1.18}
\usepgfplotslibrary{fillbetween}
\usepackage[np]{numprint}
\usepackage{listings,makecell,pas-tableur,twoopt}
\let\barre\undefined
\usepackage[pstricks]{tablvar}
\definecolor{darkgreen}{RGB}{0,100,0}
\providecommand{\red}{\color{red}}
\providecommand{\blue}{\color{blue}}
\providecommand{\green}{\color{green}}
\providecommand{\black}{\color{black}}
\providecommand{\white}{\color{white}}
\providecommand{\gray}{\color{gray}}
\providecommand{\cyan}{\color{cyan}}
\providecommand{\magenta}{\color{magenta}}
\providecommand{\yellow}{\color{yellow}}
\providecommand{\violet}{\color{violet}}
\providecommand{\orange}{\color{orange}}
\providecommand{\euro}{\texteuro}
\providecommand{\EUR}{\texteuro}
\providecommand{\ieme}{\textsuperscript{e}}
\providecommand{\degres}{\ensuremath{^\circ}}
\providecommand{\degree}{\ensuremath{^\circ}}
\providecommand{\vv}[1]{\overrightarrow{#1}}
\providecommand{\Coord}[1]{#1}
\providecommand{\decosix}{}
\providecommand{\starredbullet}{\ensuremath{\bullet}}
\pagestyle{empty}
\setlength{\parindent}{0pt}
\usepackage{adjustbox}\DeclareUnicodeCharacter{2236}{:}\DeclareUnicodeCharacter{2190}{\ensuremath{\leftarrow}}
\providecommand{\R}{\mathbb{R}}
\providecommand{\N}{\mathbb{N}}
\providecommand{\Z}{\mathbb{Z}}
\providecommand{\Q}{\mathbb{Q}}
\providecommand{\D}{\mathbb{D}}
\providecommand{\C}{\mathbb{C}}
\providecommand{\e}{{\mathrm{e}}}
\providecommand{\vect}[1]{\overrightarrow{#1}}
\providecommand{\vectt}[1]{\overrightarrow{\mathrm{#1}}}
\providecommand{\vvt}[1]{\overrightarrow{\mathrm{#1}}}
\providecommand{\np}[2][]{{#2~#1}}
\providecommand{\vv}[1]{{\overrightarrow{#1}}}
\renewcommand{\d}{{\mathrm{d}}}
\providecommand{\degres}{{^\circ}}
\providecommand{\degre}{{^\circ}}
\providecommand{\blue}{\color{blue}}
\providecommand{\red}{\color{red}}
\providecommand{\enskip}{\quad}
\providecommand{\llbracket}{\lbrack\!\lbrack}
\providecommand{\rrbracket}{\rbrack\!\rbrack}
\providecommand{\euro}{{\text{€}}}
\providecommand{\no}{{\text{n°}}}
\providecommand{\interval}[2]{{[#1;#2]}}
\providecommand{\intervFO}[2]{{[#1;#2[}}
\providecommand{\dotfill}{\ldots}

\providecommand{\ds}{\displaystyle}
\providecommand{\textding}[1]{#1}
\providecommand{\Oij}{\(\mathrm{O};\vec{\imath},\vec{\jmath}\)}
\providecommand{\Oijk}{\(\mathrm{O};\vec{\imath},\vec{\jmath},\vec{k}\)}
\providecommand{\Ouv}{\(\mathrm{O};\vec{u},\vec{v}\)}
\providecommand{\intervFF}[2]{\ensuremath{[#1;#2]}}
\renewcommand{\euro}{\ifmmode\text{\texteuro}\else\texteuro\fi}
\DeclareUnicodeCharacter{00B0}{\ifmmode^{\circ}\else\textdegree\fi}
\DeclareUnicodeCharacter{20AC}{\ifmmode\text{\texteuro}\else\texteuro\fi}
% Énoncés : Éducation nationale. Transcription et corrigés : APMEP, Denis Vergès et contributeurs. Découpage et classement : CoopMaths, Sébastien Lozano, Éric Elter et contributeurs. Adaptation de la présentation : Maths Explorer.
% CC BY-SA : https://coopmaths.fr/www/about/
% https://coopmaths.fr/alea/static/e3c/2020/tex/e3c_2020_00_sujet28_4_cor.tex
\begin{document}
\noindent\textit{APMEP / CoopMaths. CC BY-SA. Présentation adaptée par Maths Explorer.}\par\medskip
\psset{xunit=1.2cm,yunit=1.2cm,labelFontSize=\scriptstyle,showorigin=false}
	\begin{pspicture}(-1,-1.4)(10.5,3.7)
		\multido{\n=-0.2+0.2}{56}{\psline[linewidth=0.35pt,linecolor=lightgray](\n,-1)(\n,3)}
		\multido{\n=-1.0+0.2}{21}{\psline[linewidth=0.35pt,linecolor=lightgray](0,\n)(10.8,\n)}
		\multido{\n=0+1}{11}{\psline[linewidth=0.45pt](\n,-1)(\n,3)}
		\multido{\n=-1+1}{5}{\psline[linewidth=0.45pt](0,\n)(10.8,\n)}
		\psaxes[linewidth=0.95pt]{->}(0,0)(0,-1)(11,3)
		\def\Func{x 6 mul  2.71828 x exp div}
		\psplot[plotpoints=3000,linewidth=1.25pt,linecolor=red]{0}{10}{\Func}
		\uput[ur](2.5,1.5){\red $\mathcal{C}$}
	\end{pspicture}
	
	\subsection*{1.}
	La fonction \(f\) est un quotient de fonctions dérivables, le dénominateur étant non nul quel que soit le réel \(x\). On a donc sur \(\mathbb{R}\), et en particulier sur \([0\,;\,10]\) :
	\[
	f'(x) = \dfrac{6 e^x - 6x \times e^x}{(e^x)^2} = \dfrac{6 e^x (1 - x)}{e^x \times e^x} = \dfrac{6 (1 - x)}{e^x}.
	\]
	
	\subsection*{2.}
	Comme \(6 e^x > 0\) quel que soit \(x\), pour \(x \in [0\,;\,10]\), le signe de \(f'(x)\) est celui de \(1 - x\) :

		 \begin{align*}
		&1 - x > 0\\
		 \iff& 1 > x \\
		\iff& x < 1 
		\end{align*}
		La fonction \(f\) est donc croissante sur \([0\,;\,1]\), \\
		de \(f(0) = 0\) à \(f(1) = \dfrac{6 \times 1}{\e^1} = 6 \e^{-1} \approx 2,207\) ;
	
		\begin{center}
		\begin{tikzpicture}[double distance=2pt]
			\tkzTabInit{$x$/1,f'(x)/1,$f(x)$/2}{$0$,$1$, $10$}
			\tkzTabLine{,+,z, -}
			\tkzTabVar{+/$0$,-/ $6\e^{-1}$ /, +/}
		\end{tikzpicture}
	\end{center} 
	\subsection*{3.}
	Les variations de \(f\) montrent que la concentration maximale est atteinte après 1 heure et qu'elle est égale à environ \(2,2\) (mg/L).
	
	\subsection*{4.}
	Le sportif peut être contrôlé à tout moment après son injection, mais hélas il sera en infraction entre environ 43 minutes et 1h24 après l'injection car la concentration sera à ce moment supérieure à \(2\) (mg/L).

\end{document}
